\section{Introduction}\label{sec:intro}
Digital navigation systems are convenient aids for reaching unfamiliar destinations, but successful wayfinding need not imply that a traveler has learned the environment. A device may supply the next turn without requiring its user to form an integrated route representation. In a pedestrian comparison, Ishikawa \emph{et al.} reported differences in spatial learning under GPS guidance, maps, and direct experience~\cite{ishikawa_wayfinding_2008}. In a driving experiment, Ben-Elia compared a paper map with audiovisual turn-by-turn (TBT) instructions and found a landmark-recognition advantage for the map condition~\cite{ben-elia_exploratory_2021}. Evidence from habitual GPS use is also relevant, but its associations with spatial memory should not be interpreted as causal effects of an individual journey~\cite{dahmani_habitual_2020,topete_how_2024}.

\subsection{Prior Work and Research Gap}
Spatial knowledge is multidimensional. Kim and Bock found that landmark, route, and survey knowledge could develop in overlapping processes rather than in strictly separate stages~\cite{kim_landmark_2021}. This motivates testing recognition of environmental features separately from memory for turning decisions. Prior map exposure may support some forms of spatial knowledge: in a virtual-navigation experiment, Meneghetti and Pazzaglia observed advantages for map study on pointing and path drawing, but not on route retracing~\cite{meneghetti_navigating_2021}. Navigation assistance can also trade immediate wayfinding accuracy against learning of environmental configuration~\cite{munzer_navigation_2012}.

The effectiveness of navigation information depends on its representation and content. Orientation maps, the time course of digital guidance, landmark visualization, and landmark-enriched instructions have been investigated in complementary settings~\cite{lowen_spatial_2019,zhao_time_2023,cheng_effect_2022,wunderlich_landmark_2023}. Recent studies additionally consider GPS versus non-GPS conditions and learning across repeated mobile-map navigation episodes~\cite{al-srogy_spatial_2025,ying_mobile_2026}. However, findings from virtual driving, repeated navigation, or paper maps do not directly resolve a focused pedestrian comparison after \emph{one} exposure: a digital map studied \emph{before} walking without subsequent access (MAP) versus digital TBT instructions \emph{during} walking.

\subsection{Research Questions and Contributions}
We ask: \emph{RQ1}, how frequently did participants report using TBT and prior digital map study for unfamiliar routes during the last 30 days? \emph{RQ2}, how did the assigned methods differ in direction recall after one walk? \emph{RQ3}, how did they differ in landmark recognition? Prior evidence motivates directional hypotheses for both experimental outcomes: $H_{1,D}:\mu_{\mathrm{MAP},D}>\mu_{\mathrm{TBT},D}$ (RQ2) and $H_{1,L}:\mu_{\mathrm{MAP},L}>\mu_{\mathrm{TBT},L}$ (RQ3). Their null hypotheses are $H_{0,Y}:\mu_{\mathrm{MAP},Y}\leq\mu_{\mathrm{TBT},Y}$. These one-sided hypotheses were adopted as an instructor-approved correction to the earlier two-sided wording after data collection; their confirmatory status is therefore limited. The reported one-sided tests are explicitly a revised, post-data analysis, not a preregistered test.

Our contributions are: (i) a randomized, single-route comparison with an active TBT reference; (ii) separate operational measures for direction and landmark memory; (iii) a paired descriptive assessment of navigation habits, explicitly distinguished from treatment assignment; and (iv) a documented analysis pipeline with bootstrap uncertainty. Section~\ref{sec:methods} details the protocol, Section~\ref{sec:results} reports the evidence, and Section~\ref{sec:conclusion} addresses validity and implications.
